
\documentclass{article}
\usepackage{colortbl}
\usepackage{makecell}
\usepackage{multirow}
\usepackage{supertabular}

\begin{document}

\newcounter{utterance}

\centering \large Interaction Transcript for game `taboo', experiment `medium\_en', episode 18 with reflexion{-}llama{-}3.1{-}8b{-}instruct.
\vspace{24pt}

{ \footnotesize  \setcounter{utterance}{1}
\setlength{\tabcolsep}{0pt}
\begin{supertabular}{c@{$\;$}|p{.15\linewidth}@{}p{.15\linewidth}p{.15\linewidth}p{.15\linewidth}p{.15\linewidth}p{.15\linewidth}}
    \# & $\;$A & \multicolumn{4}{c}{Game Master} & $\;\:$B\\
    \hline

    \theutterance \stepcounter{utterance}  
    & & \multicolumn{4}{p{0.6\linewidth}}{
        \cellcolor[rgb]{0.9,0.9,0.9}{
            \makecell[{{p{\linewidth}}}]{
                \texttt{\tiny{[P1$\langle$GM]}}
                \texttt{You have attempted this same task before and did not succeed. These are the plans you wrote after those attempts.} \\
\texttt{Trial \#1: New plan: I used the geographical location of the target word as the clue, but I didn't provide any additional information that would help the other player narrow down the possibilities. Since the target word "georgia" is a proper noun and a state, I should have included more details that distinguish it from other states or countries. Next time, I will provide a clue that highlights a unique characteristic of the target word, such as its capital city, a notable feature, or a historical event associated with it. For example, I could have said "CLUE: Home to the 1996 Summer Olympics."} \\
\\ 
\texttt{Trial \#2: New plan: I used the geographical location of the target word as the clue, but I didn't provide any additional information that would help the other player narrow down the possibilities. Since the target word "georgia" is a proper noun and a state, I should have included more details that distinguish it from other states or countries. Next time, I will provide a clue that highlights a unique characteristic of the target word, such as its capital city, a notable feature, or a historical event associated with it. However, in this case, I also used the related word "atlanta" in my second clue, which led to the loss. Next time, I will make sure to avoid using related words in my clues.} \\
\\ 
\\ 
\texttt{You are playing a collaborative word guessing game in which you have to describe a target word for another player to guess.} \\
\\ 
\texttt{Rules:} \\
\texttt{(a) You have to reply in the form: CLUE: <some text>. Guesses from the other player will start with GUESS.} \\
\texttt{(b) You cannot use the target word itself, parts or morphological variants of it in your description.} \\
\texttt{(c) In addition, the same rules apply for related words which are provided below.} \\
\\ 
\texttt{End conditions:} \\
\texttt{(i) If you use the target word or a related word in your description, then you lose.} \\
\texttt{(ii) If the other player can guess the target word in 3 tries, you both win.} \\
\\ 
\texttt{Let us start.} \\
\\ 
\texttt{This is the target word that you need to describe and that the other player needs to guess:} \\
\\ 
\texttt{georgia} \\
\\ 
\texttt{Related words are:} \\
\\ 
\texttt{{-} state} \\
\texttt{{-} country} \\
\texttt{{-} atlanta} \\
\\ 
\texttt{Important: You are under time pressure, give short descriptions that are to the point!} \\
            }
        }
    }
    & \\ \\

    \theutterance \stepcounter{utterance}  
    & \multicolumn{4}{p{0.6\linewidth}}{
        \cellcolor[rgb]{0.9,0.9,0.9}{
            \makecell[{{p{\linewidth}}}]{
                \texttt{\tiny{[P1$\rangle$GM]}}
                \texttt{CLUE: Southern US state known for peaches and a historic civil rights movement.} \\
            }
        }
    }
    & & \\ \\

    \theutterance \stepcounter{utterance}  
    & & & \multicolumn{2}{p{0.3\linewidth}}{
        \cellcolor[rgb]{0.9,0.9,0.9}{
            \makecell[{{p{\linewidth}}}]{
                \texttt{\tiny{[GM$|$GM]}}
                \texttt{continue} \\
            }
        }
    }
    & & \\ \\

    \theutterance \stepcounter{utterance}  
    & & & \multicolumn{2}{p{0.3\linewidth}}{
        \cellcolor[rgb]{0.9,0.9,0.9}{
            \makecell[{{p{\linewidth}}}]{
                \texttt{\tiny{[GM$|$GM]}}
                \texttt{Related word 'state' (stem=state) is similar to clue word 'state' (stem=state)} \\
            }
        }
    }
    & & \\ \\

    \theutterance \stepcounter{utterance}  
    & \multicolumn{4}{p{0.6\linewidth}}{
        \cellcolor[rgb]{0.95,0.95,0.95}{
            \makecell[{{p{\linewidth}}}]{
                \texttt{\tiny{[P2$\rightarrow$P1]}}
                \texttt{Post{-}episode Reflexion (not sent to GM):} \\
\texttt{New plan: I used the geographical location of the target word as the clue, but I didn't provide any additional information that would help the other player narrow down the possibilities. Since the target word "georgia" is a proper noun and a state, I should have included more details that distinguish it from other states or countries. Next time, I will provide a clue that highlights a unique characteristic of the target word, such as its capital city, a notable feature, or a historical event associated with it. However, in this case, I also used the related word "atlanta" in my second clue, which led to the loss. Next time, I will make sure to avoid using related words in my clues. Additionally, I should have been more careful with the information I provided, as the clue "Home to the 1996 Summer Olympics" could be associated with other cities or states as well. Next time, I will try to provide a more specific and unique clue that is less likely to be associated with other words. For example, I could have said "CLUE: Known for its gold rush history and the Chattahoochee River."} \\
            }
        }
    }
    & & \\ \\

\end{supertabular}
}

\end{document}
